\documentclass{article}
\usepackage{graphicx}
\usepackage{amsmath}
\usepackage{amssymb}
\usepackage{listings}
\usepackage{xcolor}

\lstset{
  basicstyle=\ttfamily\footnotesize,
  keywordstyle=\color{blue},
  commentstyle=\color{gray},
  stringstyle=\color{red!70!black},
  language=Python,
  breaklines=true,
  frame=single,
  showstringspaces=false,
  captionpos=b,
}

\title{BSPO --- Annotated equations with verl-style pseudo code}
\author{Reading notes for plan0.md}
\date{April 2026}

\begin{document}
\maketitle

\noindent
\textbf{Purpose.} Restate the equations of \texttt{bisim\_eqs.tex} in an order
that matches the implementation, and attach a numpy-style pseudo-code block
to each variable. Variable names follow verl's \texttt{core\_algos.py} /
\texttt{compute\_policy\_loss\_*} signatures so this file doubles as an
implementation reference.

\section{Notation and verl variable map}

\begin{tabular}{l|l|l}
\textbf{Math} & \textbf{Verl name} & \textbf{Shape / semantics} \\
\hline
$\log r_{\tau,t}$ & \texttt{log\_ratio} & $(B, T)$, $B$ = num trajectories in batch \\
$\text{mask}_{\tau,t}$ & \texttt{response\_mask} & $(B, T)$, 1 for valid tokens \\
$T_\tau$ & \texttt{seq\_lens} & $(B,)$, $\texttt{response\_mask.sum(-1)}$ \\
$\hat A_\tau$ (token level) & \texttt{advantages} & $(B, T)$; GRPO already broadcasts per-trajectory scalar \\
$g$ (group) & \texttt{group\_ids} & $(B,)$, e.g. verl's \texttt{uid} field \\
$\text{dom}$ (domain) & \texttt{domain\_ids} & $(B,)$, from curriculum sampler \\
batch composition & $B = 128 \times 16$ & 128 groups, $n=16$ rollouts each (typical) \\
\end{tabular}

\medskip
\noindent
\textbf{Convention in pseudo code below}: all shapes are stated to the right
of each assignment. Masked means / sums use the standard verl pattern
\verb|(x * mask).sum(-1) / mask.sum(-1).clamp(min=1)|.

\section{Per-trajectory primitives}

\subsection{Per-token log ratio}
\[
\log r_{\tau,t} = \log \pi_\theta(a_{\tau,t}) - \log \pi_{\theta_{\text{old}}}(a_{\tau,t})
\]

\begin{lstlisting}[caption={log ratio (verl standard)}]
# (B, T) each
log_ratio = log_prob - old_log_prob                      # (B, T)
log_ratio = torch.clamp(log_ratio, min=-20.0, max=20.0)  # numerical safety
\end{lstlisting}

\subsection{Per-trajectory sign-decomposed deviation ($s_\tau^+$, $s_\tau^-$)}

\textbf{Canonical form (Q1 answer: arithmetic):}
\[
s_\tau^+ = \frac{1}{|\tau|} \sum_{t \in \tau} \max(r_t - 1, 0),
\qquad
s_\tau^- = \frac{1}{|\tau|} \sum_{t \in \tau} \max(1 - r_t, 0)
\]

Both are non-negative. $s_\tau^+ > 0$ counts tokens whose probability is
raised by the new policy; $s_\tau^-$ counts tokens whose probability is
lowered.

\begin{lstlisting}[caption={arithmetic form (per Q1 answer)}]
ratio = torch.exp(log_ratio)                                              # (B, T)
T_tau = response_mask.sum(-1).clamp(min=1)                                # (B,)

s_pos_tau = (torch.clamp(ratio - 1.0, min=0.0) * response_mask).sum(-1) / T_tau   # (B,)
s_neg_tau = (torch.clamp(1.0 - ratio, min=0.0) * response_mask).sum(-1) / T_tau   # (B,)
\end{lstlisting}

\subsection{Signed trajectory deviation}
\[
s_\tau = s_\tau^+ - s_\tau^-
\]

\begin{lstlisting}
s_tau = s_pos_tau - s_neg_tau                            # (B,)
\end{lstlisting}

\section{The two $\Delta$ operators}

\subsection{Regular two-sided clip difference $\Delta_z^{(\text{reg})}$}
\[
\Delta_z^{(\text{reg})} = \operatorname{clip}(s_z^+, \delta_z) - \operatorname{clip}(s_z^-, \delta_z)
\]
Here \verb|clip(x, delta)| means $\min(x, \delta)$ (upper-clip; $s^\pm \ge 0$
by construction). Below we use \verb|torch.clamp(x, max=delta)|.

\begin{lstlisting}
def delta_reg(s_pos: Tensor, s_neg: Tensor, delta: float) -> Tensor:
    """
    Two-sided clipped difference.
    s_pos, s_neg: (B,), delta: scalar
    returns:     (B,)
    """
    return torch.clamp(s_pos, max=delta) - torch.clamp(s_neg, max=delta)
\end{lstlisting}

\subsection{Wasserstein-style signed min-clip $\Delta_z^{(\text{W})}$}
\[
\Delta_z^{(\text{W})} = \operatorname{sgn}(s_z) \cdot
  \operatorname{clip}\bigl(\min(s_z^+, s_z^-), \, \delta_z\bigr)
\]

\begin{lstlisting}
def delta_w(s_pos: Tensor, s_neg: Tensor, delta: float) -> Tensor:
    """
    Wasserstein-style min-clip with sign from (s_pos - s_neg).
    s_pos, s_neg: (B,) in the trajectory case; (G,) for groups; (D,) for domains.
    returns: same shape
    """
    s = s_pos - s_neg                                         # (*,) signed
    clipped_min = torch.clamp(torch.minimum(s_pos, s_neg), max=delta)  # (*,)
    return torch.sign(s) * clipped_min                        # (*,)
\end{lstlisting}

\section{Advantage hierarchy (GRPO-style)}

\subsection{Trajectory advantage (shipped)}
\[
\hat A_\tau = \frac{R_\tau - \bar R_g}{\sigma_g + \varepsilon}
\]
In verl this is computed upstream (function \texttt{compute\_grpo\_outcome\_advantage}
at \texttt{core\_algos.py:267}) and broadcast per-token into the
\texttt{advantages} tensor. The shipped BSPO loss recovers a
per-trajectory scalar with:

\begin{lstlisting}
# GRPO advantage is constant along t within a trajectory; take any valid token
A_tau = (advantages * response_mask).sum(-1) / T_tau         # (B,)
\end{lstlisting}

\subsection{Group / domain advantage (NOT shipped for single-domain)}

The math definitions are retained here for reference; the implementation
intentionally does \emph{not} compute them. Per the Q13 rigorous derivation
in \S7, $\hat A_g \equiv 0$ and $\hat A_\text{dom} \equiv 0$ under verl's
GRPO full-group baseline, so the corresponding terms in Objective 2 vanish
identically. The $\lambda_{\text{Tj}} = 0$ anchor of the Objective 3 sweep
already covers Objective 2.

\[
\hat A_g = \frac{1}{n} \sum_{\tau \in g} \hat A_\tau,
\qquad
\hat A_{\text{dom}} = \frac{1}{G_d} \sum_{g : \text{dom}} \hat A_g
\]

\[
s_g^\pm = \frac{1}{n} \sum_{\tau \in g} s_\tau^\pm,
\qquad
s_\text{dom}^\pm = \frac{1}{G_d} \sum_{g \in \text{dom}} s_g^\pm
\]

If multi-domain is reintroduced, the stash pattern (\S8) shows the upstream
hooks needed to bring \texttt{group\_ids} and \texttt{domain\_ids} into the
loss function; a \texttt{scatter\_mean} helper then builds the aggregates.
The current shipped code does not touch these paths.

\section{Four objective variants}

Ordered by complexity of data dependency (simpler first).

\subsection{Objective 1 --- Simplest baseline: $\hat J_{128 \times 16}$}
\[
\hat J_{128 \times 16}
  = \widehat{\mathbb{E}}_{\tau \sim g \sim \text{dom}}
    \bigl[\Delta_\tau^{(\text{reg})} \hat A_\tau\bigr]
\]
\textbf{Uses only trajectory-level data.} No group or domain IDs needed.

\begin{lstlisting}[caption={Objective 1 --- baseline}]
# ingredients: s_pos_tau, s_neg_tau, A_tau  (all (B,))
delta_tau_reg = delta_reg(s_pos_tau, s_neg_tau, config.bspo_delta_tau)   # (B,)
per_traj = delta_tau_reg * A_tau                                         # (B,)

# Broadcast to per-token to reuse verl's agg_loss:
per_token_obj = per_traj.unsqueeze(-1).expand_as(advantages)             # (B, T)
loss = -agg_loss(per_token_obj, response_mask, loss_agg_mode)
\end{lstlisting}

\subsection{Objective 4 --- Wasserstein Penalty only}
\[
\hat J = \widehat{\mathbb{E}}
  \Bigl[ s_\tau \hat A_\tau
       - \lambda_{\text{Tj}} \max\Bigl(0,\, \tfrac{\min(s_\tau^+, s_\tau^-)}{\delta_\tau} - 1\Bigr) \Bigr]
\]
\textbf{Trajectory-level only.} Uses the \emph{raw} signed deviation
$s_\tau$ (no clip on the advantage term) plus a soft penalty that activates
when $\min(s^+, s^-) > \delta_\tau$.

\begin{lstlisting}[caption={Objective 4 --- penalty only}]
penalty = torch.clamp(
    torch.minimum(s_pos_tau, s_neg_tau) / config.bspo_delta_tau - 1.0,
    min=0.0,
)                                                                        # (B,)
per_traj = s_tau * A_tau - config.bspo_lambda_tj * penalty               # (B,)

per_token_obj = per_traj.unsqueeze(-1).expand_as(advantages)             # (B, T)
loss = -agg_loss(per_token_obj, response_mask, loss_agg_mode)
\end{lstlisting}

\subsection{Objective 3 --- Wasserstein Penalty (clipped $\Delta^{(W)}$ + penalty)}
\[
\hat J = \widehat{\mathbb{E}}
  \Bigl[ \Delta_\tau^{(\text{W})} \hat A_\tau
       - \lambda_{\text{Tj}} \max\Bigl(0,\, \tfrac{\min(s_\tau^+, s_\tau^-)}{\delta_\tau} - 1\Bigr) \Bigr]
\]
\textbf{Trajectory-level only.} Clipped advantage term plus the same soft
penalty as Objective 4.

\begin{lstlisting}[caption={Objective 3 --- W + penalty, traj only}]
delta_tau_w = delta_w(s_pos_tau, s_neg_tau, config.bspo_delta_tau)       # (B,)
penalty = torch.clamp(
    torch.minimum(s_pos_tau, s_neg_tau) / config.bspo_delta_tau - 1.0,
    min=0.0,
)                                                                        # (B,)
per_traj = delta_tau_w * A_tau - config.bspo_lambda_tj * penalty         # (B,)

per_token_obj = per_traj.unsqueeze(-1).expand_as(advantages)             # (B, T)
loss = -agg_loss(per_token_obj, response_mask, loss_agg_mode)
\end{lstlisting}

\subsection{Objective 2 --- Strict Wasserstein (NOT implemented; reduces to Obj 3 at $\lambda=0$)}
\begin{align*}
\hat J =\ & \widehat{\mathbb{E}}_{\tau \sim g \sim \text{dom}}
  \bigl[\Delta_{\text{dom}}^{(\text{W})} \hat A_{\text{dom}}
      + \Delta_g^{(\text{W})} \hat A_g
      + \Delta_\tau^{(\text{W})} \hat A_\tau\bigr]
\end{align*}

\textbf{Why not implemented.} Under verl's GRPO full-group baseline,
$\hat A_g \equiv 0$ and $\hat A_\text{dom} \equiv 0$ (Q13 derivation, \S7).
The strict-Wasserstein objective therefore collapses to
$\widehat{\mathbb{E}}_\tau [\Delta_\tau^{(W)} \hat A_\tau]$, which is
exactly Objective 3 with $\lambda_{\text{Tj}} = 0$. The sweep plan
includes $\lambda = 0$ as an anchor, so Obj 2's gradient signal is
covered for free without any 3-level aggregation code.

No pseudo code is shipped for this objective in the single-domain
round. If multi-domain is reintroduced, the 3-level aggregation
template appears in \S4.2 (group/domain advantage) and \S8 (upstream
plumbing hooks from the stash).

\section{Reference block: motivation only (not implemented)}

\textbf{Q2 answer}: the ${}^{(\text{GSPO})}\tilde\delta_\tau$ expression
below has \emph{nothing to do with the GSPO loss}. It was written in the
same visual block purely as a reminder that $\delta$-hparams here should
be \emph{comparable in magnitude} to GSPO's. The four objectives do not
use it.

The definition of $\hat d(z)$ is the exact implementation of $\hat d(\text{dom})$,
$\hat d(g)$, etc.: the mean-absolute-deviation of rewards for
trajectories drawn from the distribution $\mathcal{P}(z)$.

\textbf{Single-domain simplification}: we are only implementing the
single-domain case for now. Any term that is constant across a
single-domain batch is dropped. In particular:

\begin{itemize}
\item $\hat d(\text{dom})$, $R_{\text{max res}}(\text{dom})$, and
  $\delta_D$ collapse to scalars constant over the batch and drop from
  any gradient.
\item Only $\hat d(g)$ and per-trajectory terms retain per-sample
  variation, but since $\hat d$ itself does not appear in the 4
  objectives, we do not need to compute it at all in the loss function.
\end{itemize}

\begin{align*}
\hat d(z) &= \mathbb{E}_{\tau \sim \mathcal{P}(z)}\bigl[ |R(\tau) - \mathbb{E}_{\tau'\sim\mathcal{P}(z)}[R]| \bigr]
             \quad \text{(definition only --- not in any implemented loss)}
\end{align*}

No pseudo code for this section --- nothing here enters the gradient
path of Objectives 1--4 in the single-domain implementation.

\section{Single-domain consequence for Objective 2 (rigorous derivation)}

\textbf{Claim.} Under verl's GRPO advantage (full-group mean baseline,
full-group std) and single-domain batches, Objective 2 (Strict Wasserstein)
is \emph{identically equal} to Objective 3 with $\lambda_{\text{Tj}} = 0$.

\textbf{Verl's GRPO advantage} (from
\texttt{verl/trainer/ppo/core\_algos.py:267-326}, function
\texttt{compute\_grpo\_outcome\_advantage}):

\begin{align}
R_\tau &:= \sum_{t} \text{token\_level\_reward}_{\tau,t} \quad \text{(scalar per trajectory)} \\
\bar R_g &:= \frac{1}{n}\sum_{\tau \in g} R_\tau \quad \text{(\texttt{torch.mean}, full group)} \\
\sigma_g &:= \operatorname{std}(\{R_\tau : \tau \in g\}) \quad \text{(\texttt{torch.std}, full group)} \\
\hat A_\tau &:= \frac{R_\tau - \bar R_g}{\sigma_g + \varepsilon} \quad \text{(broadcast to every token in \(\tau\))}
\end{align}

\textbf{Derivation.} By definition of $\bar R_g$:
\begin{align}
\sum_{\tau \in g} (R_\tau - \bar R_g)
&= \sum_{\tau \in g} R_\tau - n \bar R_g \\
&= n \bar R_g - n \bar R_g = 0 \quad \text{(exactly)}
\end{align}

Therefore:
\begin{align}
\sum_{\tau \in g} \hat A_\tau
&= \frac{1}{\sigma_g + \varepsilon} \sum_{\tau \in g} (R_\tau - \bar R_g)
 = \frac{0}{\sigma_g + \varepsilon} = 0 \quad \text{(exactly)}
\end{align}

Consequently:
\begin{align}
\hat A_g     &:= \frac{1}{n} \sum_{\tau \in g} \hat A_\tau \equiv 0 \\
\hat A_{\text{dom}} &:= \frac{1}{G_d} \sum_{g \in \text{dom}} \hat A_g \equiv 0
\end{align}

\textbf{Substitute into Objective 2:}
\begin{align}
\hat J_{\text{strict W}}
&= \widehat{\mathbb{E}}_\tau \bigl[\Delta_{\text{dom}}^{(W)} \hat A_{\text{dom}}
        + \Delta_g^{(W)} \hat A_g
        + \Delta_\tau^{(W)} \hat A_\tau \bigr] \\
&= \widehat{\mathbb{E}}_\tau \bigl[\Delta_{\text{dom}}^{(W)} \cdot 0
        + \Delta_g^{(W)} \cdot 0
        + \Delta_\tau^{(W)} \hat A_\tau \bigr] \\
&= \widehat{\mathbb{E}}_\tau \bigl[\Delta_\tau^{(W)} \hat A_\tau \bigr]
\end{align}

Objective 3 at $\lambda_{\text{Tj}} = 0$:
\begin{align}
\hat J_{\text{W penalty}}\bigm|_{\lambda=0}
&= \widehat{\mathbb{E}}_\tau \bigl[\Delta_\tau^{(W)} \hat A_\tau - 0\bigr]
 = \widehat{\mathbb{E}}_\tau \bigl[\Delta_\tau^{(W)} \hat A_\tau \bigr]
\end{align}

\textbf{Conclusion.} $\hat J_{\text{strict W}} \equiv \hat J_{\text{W penalty}}\bigm|_{\lambda=0}$ as an identity (not an approximation), for any single-domain batch under verl's GRPO advantage with full-group mean baseline.

\textbf{Implication for implementation.} Under the "math always wins" rule:
Objective 2 is redundant in the single-domain case --- any run we do with
$\lambda_{\text{Tj}} = 0$ in Objective 3 already covers Objective 2
exactly. The effective ablation set is:
\[
\{\text{Obj 1}, \ \text{Obj 3 (with }\lambda\in\{0, 0.001, 0.005\}\text{)}, \ \text{Obj 4 (with }\lambda\in\{0, 0.001, 0.005\}\text{)}\}
\]
We do not need to implement the group / domain aggregation machinery at
all for the first ablation round; the $\lambda=0$ point of Obj 3 already
gives the "Strict Wasserstein" result. Group-level plumbing can be added
later when multi-domain is reintroduced.

\textbf{When the reduction breaks.} If we change the baseline (e.g. to
leave-one-out $\bar R_g^{-\tau} = \frac{1}{n-1}\sum_{\tau' \ne \tau} R_{\tau'}$)
then $\sum_{\tau \in g}(R_\tau - \bar R_g^{-\tau}) \ne 0$ in general, and
$\hat A_g$ becomes non-zero. For now we stay on verl's default full-group
baseline.

\section{Shipped signature (single-domain, Obj 1/3/4)}

The code that actually ships (\texttt{verl/verl/trainer/ppo/core\_algos.py}):

\begin{lstlisting}[caption={BSPO loss signature as shipped}]
@register_policy_loss("bspo")
def compute_policy_loss_bspo(
    old_log_prob: torch.Tensor,         # (B, T)
    log_prob:      torch.Tensor,         # (B, T)
    advantages:    torch.Tensor,         # (B, T), GRPO-normalized constant-along-t
    response_mask: torch.Tensor,         # (B, T)
    loss_agg_mode: str = "seq-mean-token-mean",   # ignored; hard-coded inside
    config:        Optional[ActorConfig] = None,
    rollout_is_weights: torch.Tensor | None = None,   # ignored (Q9)
) -> tuple[torch.Tensor, dict[str, Any]]:
    ...
\end{lstlisting}

\textbf{No extra kwargs.} Since Obj 2 reduces to Obj 3 at $\lambda=0$ under
single-domain + GRPO (Q13), we do not need \texttt{trajectory\_rewards}
or \texttt{group\_ids} on the signature, nor the matching plumbing in
\texttt{dp\_actor.py} / \texttt{megatron\_actor.py}. Zero changes to the
actor workers.

\subsection{Future multi-domain hook (reference only)}

If multi-domain is reintroduced later, the stash pattern at
\texttt{verl}@\texttt{stash@\{0\}} is the template: add the two kwargs to
the signature, and add a small hook in \texttt{dp\_actor.py} that extracts
\texttt{token\_level\_rewards} + \texttt{uid} from \texttt{model\_inputs}
and passes them via \texttt{**extra\_kwargs}. \texttt{megatron\_actor.py}
needs a mirror hook with a prompt-hash fallback for the case where
micro-batching drops the \texttt{non\_tensor\_batch} fields. None of this
is shipped now.

\section{Loss-aggregation math (Q7)}

The \texttt{bisim\_eqs.tex} objective is written as
$\widehat{\mathbb{E}}_{\tau \sim g \sim \text{dom}}[\,\cdot\,]$ --- an
expectation over trajectories in the empirical batch. Concretely:
\[
\widehat{\mathbb{E}}_\tau[f(\tau)] = \frac{1}{B} \sum_{\tau = 1}^{B} f(\tau)
\]
Each trajectory contributes equally, regardless of its token count $T_\tau$.

For verl's \texttt{agg\_loss(per\_token\_obj, mask, mode)}:
\begin{itemize}
\item \texttt{token-mean}: $\frac{\sum_{\tau,t} L_{\tau,t} \cdot \text{mask}}{\sum \text{mask}} = \frac{\sum_\tau T_\tau \cdot f(\tau)}{\sum_\tau T_\tau}$. \textbf{Length-weighted.} Does not match $\widehat{\mathbb{E}}_\tau$ in general.
\item \texttt{seq-mean-token-mean}: $\frac{1}{B}\sum_\tau \left[\frac{\sum_t L_{\tau,t} \cdot \text{mask}}{T_\tau}\right] = \frac{1}{B}\sum_\tau f(\tau)$. \textbf{Matches.}
\item \texttt{compute scalar directly}: mathematically identical to the seq-mean form above, just skips \texttt{agg\_loss}.
\end{itemize}

\textbf{Verdict (per user's "math always wins" rule):} only
\texttt{seq-mean-token-mean} is correct. No sweep needed on this axis. BSPO
will hard-code \texttt{loss\_agg\_mode = "seq-mean-token-mean"} the same way
FACPO does.

\section{Hyperparameter list (as shipped)}

\begin{tabular}{l|l|l}
\textbf{Name} & \textbf{Default (yaml)} & \textbf{Appears in} \\
\hline
\texttt{bspo\_variant}    & \texttt{simplest} & enum \{\texttt{simplest}, \texttt{w\_penalty}, \texttt{w\_penalty\_only}\}; selects Obj 1 / 3 / 4 \\
\texttt{bspo\_delta}      & $3 \times 10^{-4}$ & scalar clip threshold at GSPO scale (not GRPO's $\sim 0.2$) \\
\texttt{bspo\_lambda\_tj} & $1 \times 10^{-3}$ & soft penalty coeff; used by Obj 3 + Obj 4 only \\
\end{tabular}

Three fields live both in \texttt{verl/trainer/config/actor/actor.yaml}
(YAML defaults, per plan0.md's "yaml > bash" preference) and as typed
fields on \texttt{ActorConfig} in \texttt{verl/workers/config/actor.py}.
Selected at submit time via
\texttt{actor\_rollout\_ref.actor.bspo\_variant=...} Hydra overrides in
\texttt{bisimpo/submit\_bspo.sh}.

Fixed choices inside the loss function (not knobs):
\begin{itemize}
\item $s^\pm$ form: arithmetic \texttt{max(ratio $-$ 1, 0)} / \texttt{max(1 $-$ ratio, 0)} (Q1).
\item \texttt{loss\_agg\_mode}: hard-coded to \texttt{seq-mean-token-mean} (Q7 math verdict).
\item \texttt{sign(s\_tau).detach()} on the Wasserstein sign (Q8).
\item \texttt{rollout\_is\_weights} accepted but never read (Q9).
\end{itemize}

Why $3 \times 10^{-4}$ for \texttt{bspo\_delta}: GSPO paper uses
$\epsilon \approx 3 \times 10^{-4}$ for a log-ratio clip. BSPO's $s^\pm$
uses arithmetic $(r - 1)$, linearly equivalent to $\log r$ near $r = 1$,
so the GSPO magnitude is a reasonable anchor. Phase 2 will sweep around
this value.

\section{Ablation plan (as executed)}

\textbf{Phase 1 --- identify winning objective.} Three runs at default
hparams:
\begin{itemize}
\item \texttt{cbsp101}: \texttt{bspo\_variant=simplest}, $\delta=3 \times 10^{-4}$ (Obj 1)
\item \texttt{cbsp103}: \texttt{bspo\_variant=w\_penalty}, $\delta=3 \times 10^{-4}$, $\lambda=10^{-3}$ (Obj 3)
\item \texttt{cbsp104}: \texttt{bspo\_variant=w\_penalty\_only}, $\delta=3 \times 10^{-4}$, $\lambda=10^{-3}$ (Obj 4)
\end{itemize}

120 training steps each, \texttt{test\_freq=10}, \texttt{val\_max\_samples=128},
C11 compute profile (4 nodes, no offload, \texttt{gmu=0.4}).
Wall clock: $\approx 18$ h per run, three run concurrently on 12 of 31
compute nodes.

(Obj 2 not run; covered by $\lambda=0$ anchor in Phase 2 of Obj 3.)

\textbf{Phase 2 --- sweep hparams of the Phase 1 winner.}

\begin{itemize}
\item If winner is Obj 1: sweep $\delta \in \{1 \times 10^{-4},\, 3 \times 10^{-4},\, 1 \times 10^{-3},\, 3 \times 10^{-3}\}$. 4 additional runs (default already covered).
\item If winner is Obj 3 / Obj 4: sweep
  $(\delta, \lambda) \in \{10^{-4}, 3 \times 10^{-4}, 10^{-3}\} \times \{0,\, 10^{-3},\, 5 \times 10^{-3}\}$.
  9 runs total. The $\lambda = 0$ anchor exactly reproduces the dropped
  Objective 2.
\end{itemize}

\section{Shipped files --- traceability}

\begin{tabular}{l|l}
\textbf{File} & \textbf{Change} \\
\hline
\texttt{verl/trainer/ppo/core\_algos.py} & $+$167 lines: \texttt{compute\_policy\_loss\_bspo} + metrics + registration \\
\texttt{verl/trainer/config/actor/actor.yaml} & $+$10 lines: yaml defaults for 3 BSPO fields \\
\texttt{verl/workers/config/actor.py} & $+$4 lines: dataclass fields \\
\texttt{verl/my\_scripts/k8s/config/combo\_40Bra.yaml} & $+$32 lines: \texttt{cdbgbspo}, \texttt{cbsp101/103/104} \\
\texttt{bisimpo/submit\_bspo.sh} & new file: combo $\to$ Hydra overrides + delegates to production \texttt{submit.sh} \\
\texttt{bisimpo/launch\_phase1.sh} & new file: submits \texttt{cbsp101/103/104} as a batch \\
\texttt{iter\_kuberay\_32nodes\_verl\_training/submit/submit.sh} & relaxed combo-id regex (1 line) \\
\end{tabular}

Not touched, deliberately: \texttt{dp\_actor.py}, \texttt{megatron\_actor.py}
(see \S8 for why).

\end{document}
